\chapter{Introduzione}
\label{cha:introduction}

\par
RDF (Resource Description Framework) è la specifica definita da W3C per la rappresentazione di ontologie e grafi di conoscenze
nell'ambito del web semantico. È una specifica essenziale con il potenziale per essere estesa,
come ad esempio succede con OWL (Web Ontology Language), una sua potente estensione utilizzata per definire
ontologie complete. Di base è stata creata per la condivisione di risorse web, ognuna identificata da un IRI univoco
teoricamente accessibile in rete, ha dato il via alla creazione di un esteso ecosistema di applicativi e librerie
software per l'interrogazione e la manipolazione di grafi RDF. Nella pratica, molti DBMS sono stati integrati o creati con
il supporto per interrogazioni SPARQL (SPARQL Protocol and RDF Query Language), il protocollo e linguaggio standardizzato
sempre da W3C per interrogare basi di conoscenza nel web semantico.
\par
Anche se il semantic web non ha raggiunto l'adozione e l'impatto sperato al momento della creazione di questi standard,
rimane comunque alla base di uno dei più vasti e avanzati sistemi per la manipolazione di basi di conoscenza. Inoltre,
grazie alla sua estensibilità by-design esso permette facilmente di valicare i domini della rete per essere utilizzato
in ambienti fondamentalmente diversi rispetto all'originale idea del semantic web.
Tramite la creazione di ontologie specifiche per dei dati domini è possibile rappresentare qualsiasi tipo di dato
tramite modelli RDF, mantenendo la stessa espressività del dato originale e quindi la completezza.
I vantaggi con questo approccio sono molteplici, in primis si andrebbe a creare una rappresentazione estensibile per il
dominio preso in analisi, permettendo a chiunque abbia accesso all'ontologia di espandere le possibilità di rappresentazione
del modello senza compromettere l'interoperabilità con altri sistemi non aggiornati. Questo non è invece possibile quasi mai
con le rappresentazioni native solitamente utilizzate per i loro domini, focalizzate più sulle prestazioni che sull'estensibilità
(si veda il capitolo relativo a 3Dont).
Inoltre, con questo approccio, si può andare a sfruttare la potenza dei sistemi di ragionamento sui modelli utilizzati,
andando a definire regole per espandere la conoscenza nelle nostre rappresentazioni e potenzialmente anche ottenere un guadagno
dal punto di vista della memoria di massa utilizzata, se il motore che stiamo andando ad utilizzare supporta il ragionamento ``virtuale''
a tempo di interrogazione (come nel caso di graphdb o owlready2, analizzati nel capitolo~\ref{cha:benchmark}).
\par
In questa tesi si andrà quindi ad analizzare la fattibilità tecnica di questo approccio, con una particolare focalizzazione sull'ambito
dei big data, dati che se rappresentati come grafi di conoscenza superano le decine di milioni di nodi.
Queste grandi quantità di dati pongono appunto un problema all'ecosistema utilizzato. Se il problema della rappresentazione è per
così dire ``risolto'' dalla grande espressività ed estensibilità by-design di RDF, quello delle prestazioni è di altra natura.
Se vogliamo utilizzare RDF come modello di salvataggio per big data in modo generico dobbiamo fare i conti con dei pattern di accesso
fondamentalmente diversi da quelli tipici del semantic web. Si andrà spesso ad analizzare l'interezza del modello oppure
delle sottoporzioni per andare a condurre delle analisi su questi dati, o produrre delle rappresentazioni in formati alternativi
(grafici, ad esempio) per cui è necessario accedere in lettura a grandi porzioni di dati. Stiamo quindi parlando di query di tipo OLAP
contro il tipico carico di lavoro OLTP del semantic web: la sua natura predispone la ricerca per singole risorse ed elementi collegati,
facilmente ottimizzabile con l'utilizzo di indici a vario livello. La maggior parte dei DBMS con supporto per SPARQL non sono quindi
ottimizzati per query che contengono scansioni su grandi porzioni del modello quando sull'accesso casuale e sul collegamento dei nodi.
Questa è appunto la sfida più grande all'utilizzo diffuso dell'approccio proposto, e in questa tesi si andranno quindi soprattutto
ad analizzare le prestazioni di un certo numero di motori di query SPARQL nell'esecuzione di carichi di lavoro specifici,
con lo scopo di valutarne la predisposizione all'utilizzo con big data di dominio diverso dal web semantico.
Nello specifico si andranno ad impiegare datasets e query utilizzate dal framework 3Dont, approfondito nel capitolo dedicato,
all'interno di vari motori la cui prerogativa principale è che siano eseguibili localmente. Questo è importante soprattutto per
poter confrontare le prestazioni dei vari motori nello stesso ambiente ma anche perché se si volessero poi utilizzare come base
per l'accesso ai modelli RDF è importante che possano essere integrati all'interno degli applicativi che li utilizzerebbero.

\section{Lavori correlati}
\label{sec:prior_works}

L'ambiente di ricerca e sviluppo che negli anni si è creato intorno a RDF e al semantic web è ricco e pieno di spunti.
Vale quindi la pena analizzare alcuni lavori correlati a questa tesi per il contenuto e temi trattati come per l'approccio
alternativo allo standard.
Si riporta quindi una breve lista di ricerche e progetti che hanno usato RDF appunto come formato per l'organizzazione 
di dati, sfruttando le sue capacità aggregative, oltre che benchmark preesistenti, con una comparazione negli aspetti
qui proposti e nelle loro differenze.

\subsection{Progetti rilevanti}

Un progetto molto interessante è Linked Building Data, che aggrega dati da diverse sorgenti per monitorare lo sviluppo
dei cantieri. Parte da una rappresentazione RDF dei dati BIM e usa l'ontoligia V4Design come base di costruzione semantica.
Integra inoltre GEOSPARQL per la correlazione di dati geospaziali insieme all'ontologia FOG per la rappresentazione dei punti.
Il punto chiave del progetto poi è il collegamento con la sensoristica esistente, è un problema comune (perlaltro affrontato 
anche da 3Dont) quello dell'integrazione di dati da sorgenti così disparate, mantendendo una coesione spaziotemporale,
basti pensare al gran numero di formati proprietari usati dai sensori sul mercato, le diverse proiezioni geografiche e
temporali (ad esempio i sensori GNSS spesso usano il tempo GPS, leggermente ma significativamente scostato da quello UTC).
Un aspetto interessante sicuramente il fatto che vengano usate molte ongologie già esistenti per la definizione dei dati,
rendendo ancora più semplice l'interoperabilità, anche se GEOSPARQL limita l'utilizzo ai soli triple store che lo supportano,
come GraphDB e Virtuoso.

\par Nell'albito del progetto WaterFRAME inoltre è stato sperimentato l'utilizzo di LDES (Linked Data Event Streams), uno
standard basato su RDF per la gestione di eventi sequenziali e serie di dati temporali. L'utilizzo di Time Series Snippets 
permette di ridurre la quantità di dati da processare, andando a compattare senza perdida di conoscenza le triple che descrivono
una serie temporale in un solo valore letterale.
Viene poi usato Virtuoso, pratico data la sua integrazione con GEOSPARQL (che anche qui viene utilizzato), che garantisce le
prestazioni necessarie.

\par Anche Brick prova a risolvere il problema dell'integrazione sensoristica, diventato chiave per l'intepretazione
e l'ottimizzazione dei consumi energetici negli edifici. Usando appunto un'ontologia comune sarebbe possibile avere software
che funzionano per un vasto numero di edifici senza nessuna modifica, anche se la completezza nella rappresentazione
è una sfida complidata. Il risultato è un framework open source (https://brickschema.org/) estensibile by-default. 
Sistemi come questi sono essenziali e condividono l'idea alla base del semantic web, evitano il problema del vendor lock-in
e in questi casi raggiungono anche delle prestazioni ottime.

\par Questi progetti rendono l'idea di come sia possibile effettivamente usare RDF per un approccio più aperto alla condivisione
dei dati, ma non tengono in considerazione il problema delle prestazione. La quantità di dati utilizzata raramente riesce
a raggiungere il milione di triple quando convertita in RDF pertanto non pone nessun particolare problema alla maggior parte
dei motori SPARQL più popolari. Ma per questo esistono diversi benchmark, analizzati nella prossima sezione.

\subsection{Benchmark esistenti}

Naturalemente esiste un gran numero di insiemi di datasets e query che possiamo usare per misurare le prestazioni dei vari
triple store qui considerati, ma nessuno sembrerebbe rispecchiare in modo preciso e complete il carico di lavoro che si andrebbe 
ad incontrare se si utilizzasse RDF come formato generico di rappresentazione. Se pensiamo agli accessi che normalmente vengono
fatti su un file, la maggior parte sono letture di grandi porzioni dello stesso. Questo cambierebbe forse avendo a supporto
le query SPARQL ma dall'esperienza di 3DOnt il grande collo di bottiglia sono appunto gli enormi risultati e il poco vantaggio
dato dall'indicizzazione classica nelle grandi scansioni. 

\par BowlognaBench ad esempio parte sì da una base di dati consistente (i registri accademici dell'Università di Bologna) ma
nel vasto numero di query analitiche definite la dimensione dei risultati è sempre molto ridotta. È quindi comunque interessante
ma incompleto per il nostro caso d'uso. Soltanto la quarta query avrebbe un insieme di risultati largo ma prendendo in considerazione
dati privati degli studenti non è pubblicamente localmente.
\par Stesso vale per SSB-QB4OLAP, il cui volume dei risultati è soltanto una frazione del volume totale dei dati, nonostante
stiamo sempre parlando di query OLAP. É comunque interessante confrontare i risultati qui ottenuti con quelli di questi 
altri benchmark, si vanno a rendere chiare le differenze di scopo nelle implementazioni dei vari motori.
\par LDBC Semantic Publishing Benchmark invece si mostra interessante su due delle query definite. La quarta infatti permette
di limitare le dimensioni del risultato, scegliento un valore di limite sufficientemente elevato si possono analizzare appunto
le prestazioni dei triple store in più scansioni complete dei dati. Stesso discorso per la seconda query nella categoria avanzata,
simile a quella appena citata non ha nessun limite nella dimensione dei risultati.
Il benchmark in realtà è molto più complesso di così e andrebbe a simulare il lavoro su una base dati condivisa con operazioni di
modifica e aggregazione concorrenti. Il nostro caso d'uso invece richiede soltanto un accesso sequenziale da parte di un solo utente,
per cui le tecnologie implementate per aumentare il parallelismo dei triple store non forniscono alcun guadagno.
\par Va citato inoltre Smart City RDF Benchmark, che usa dati geospaziali proveniente dai sensori della città di Firenze, che però
richiede complesse capacità di indicizzazione del testo e accede solo a ridotte porzioni temporali per ogni query, trattandosi
di dati storici urbani.


\section{3Dont}
\label{sec:3dont}

\subsection{Ambito}

La nuvola di punti è un formato molto diffuso per la rappresentazione di dati tridimensionali geospaziali. Questo formato viene usato
da ambiti che vanno dall'urbanistica ai beni culturali e alle scienze forestali. La maggior parte dei processi di scansione e 
acquisizione di ambienti tridimensionali produce come primo risultato una nuvola di punti, su cui possono essere poi eseguiti un
vasto numero di processi, dalla segmentazione istanze all'analisi numerica delle sue caratteristiche. Una singola nuvola di punti
può rappresentare un piccolo oggetto come un'intera città, e a variare della risoluzione e della dimensione può arrivare ad avere
centinaia di milioni di punti. 
\par Esistono diversi formati per la condivisione di nuvole di punti, come PCD (Point Cloud Data), LAS (LASer) e LAZ (LASzip),
che in pratica permettono di definire una serie di campi presenti per ogni punto e la lista di punti della nuvola. Questi formati
sono più o meno complessi, LAS ad esempio permette solo 11 formati di nuvole di punti diverse (come con o senza tempo GPS, colore, etc.)
anche se grazie a questa caratteristica LASzip permettere di comprimere file LAS senza perdita di dettaglio in modo molto efficace.

\subsection{Il framework}

3Dont è un framework basato su RDF che permette di modellare nuvole di punti attraverso un'ontologia definita ma estensibile.
In fase di creazione del grafo classifica i punti e gli oggetti della nuvola in base all'ambito scelto (urbano o beni culturali)
secondo una struttura gerarchica, per poi eseguire una fase di ragionamento che materializza delle informazioni aggiuntive come la
correlazione spaziale, classificazione tassonomica, proprietà dei materiali e così via.
\par Il modello così creato è pronto per essere visualizzato ed esplorato, anche attraverso delle interrogazioni in linguaggio naturale,
che grazie a un LLM e un'analisi dell'ontologia vengono tradotte in linguaggio SPARQL ed eseguite. Questo permette anche a individui
senza una formazione specifica su nuvole di punti, ontologie o informatica (come ad esempio un pianificatore urbanistico) 
di interagire con i dati e supportare il proprio lavoro in autonomia.

\subsection{Vantaggi}

Come spiegato nell'introduzione, uno dei vantaggi principali di RDF rispetto ai corrispondenti formati nativi è l'estensibilità.
3Dont specifica un'ontologia base con un vasto numero di elementi già definiti per i vari ambiti a cui si rivolge (classificazione
delle superfici erbose e della massa verde, componenti architettonici, arredo urbano etc.) ma non mira ad essere completa, un utente
potrebbe senza sforzo aggiungere nuove classi, ad esempio nuovi materiali con nuove caratteristiche, senza compromettere l'integrazione
con la strumentazione esistente.
\par Un secondo vantaggio lo troviamo nell'espressività che di base il framework riesce a dare. Se partiamo da formati come LAS che
permettono di avere solo un numero limitato di caratteristiche per ogni punto, con 3Dont il potenziale è infinito; non dobbiamo
preoccuparci di problemi come l'assegnazione di più etichette di classificazione a un singolo punto ma invece possiamo integrare i
dati con informazioni aggiuntive come le relazioni tra i punti, tra macro entità (insiemi di punti), o informazioni spaziotemporali.
\par Nessuna di queste ultime caratteristiche è supportata nei formati di nuvole di punti classiche e l'utilizzo di un solo formato
di supporto per questa vasta gamma di informazioni è senza dubbio un ulteriore vantaggio del framework.
\par Ma uno dei punti di forza maggiore di 3Dont è la possibilità di interrogare i dati con un linguaggio potente come SPARQL,
in questo modo possiamo accedere in modo casuale alle nuvole di punti, esprimere potenti filtri e condurre analisi numeriche
senza bisogno di iterare sull'interezza dei punti.

\subsection{Visualizzatore}

Il framework mette disposizione, oltre allo strumento di creazione dei grafi popolati dai punti, un visualizzatore per questi grafi.
Caricando l'ontologia popolata in uno dei vari formati disponibili per rdf (turtle, n-triples, rdf-xml etc.) permette di creare
un progetto persistente con un'interfaccia diretta alla nuvola di punti, che può essere esplorata astraendo sopra alle nozioni di
grafo e SPARQL, interagendo con i punti usando mouse a tastiera e interrogazioni in linguaggio naturale. 
\par Il visualizzatore permette di visualizzare fondamentalmente quattro categorie diverse di risultati:
\begin{itemize}
\item Query di selezione: vengono evidenziati i punti ritornati dalla query, come ad esempio quelli appartenenti a un dato
oggetto oppure tutti i punti di un certo materiale.
\item Query scalari: i punti vengono colorati secondo una scala colorata lineare in base a un valore che viene ritornato dalla
query, come ad esempio la temperatura in una data acquisizione oppure una certa caratteristica del materiale dei punti.
\item Query di classificazione: tutti i punti appartenenti a un certo oggetto vengono colorati allo stesso modo, in modo univoco.
Si può così visualizzare la nuvola completamente segmentata, semanticamente o per istanza, a qualunque punto desiderato della scala
tassonomico o mereologica: un singolo punto può essere contemporaneamente parte di un capitello, della colonna a cui appartiene il capitello,
della facciata a cui appartiene la colonna e l'edificio a cui appartiene la facciata; è disponibile quindi una segmentazione per
diversi livelli di dettaglio.
\item Query tabellari: il risultato non è visualizzato graficamente attraverso la nuvola di punti ma in una tabella a lato, è il caso
delle query che ritornano singoli valori numerici o comunque valori che non sono collegati a singoli punti della nuvola.
\end{itemize}

\subsection{Implementazione}

La sfida principale nell'implementazione di 3Dont è sicuramente il raggiungimento di prestazioni compatibili con l'interattività
di un'interfaccia grafica e di un applicativo destinato a un'utenza inesperta. 
Nuvole di punti di medie dimensione (meno di 5 milioni di punti) possono richiedere per la query iniziale, che domanda coordinate e 
colori di ogni punto, diversi minuti di esecuzione, ma si può scendere a meno di una decina di secondi con qualche accorgimento.
\par Il principale fattore determinante (oltre ovviamente all'hardware utilizzato, che però ovviamente non dipende dall'implementazione
del framework) è il motore di query utilizzato e la sua configurazione. Nello specifico viene impegato QLever (TODO link) con vocabolario
compresso su disco e i prefissi dell'ontologia codificati. Viene quindi usata la versione integrata di QLever, disponibile come libreria
c++, che permette di ridurre al minimo l'utilizzo di risorse dato dall'IPC (Inter Process Communication) e dalla deserializzazione
dei risultati, che invece sarebbe necessaria usando SPARQL tramite l'interfaccia HTTP.

TODO FOTO
